\documentclass[11pt, a4paper]{article}

\usepackage[UTF8]{ctex}
\usepackage{geometry}
\usepackage{titlesec}
\usepackage{enumitem}
\usepackage{setspace}
\usepackage{hyperref}
\usepackage{xcolor}

\geometry{a4paper, margin=2.5cm}
\setstretch{1.25}

\hypersetup{
	colorlinks=true,
	linkcolor=black,
	urlcolor=blue
}

\title{剧本企划书}
\author{黄浩宇}
\date{\today}

\begin{document}
	
	\maketitle
	\newpage
	
	\section{【项目概述】}
	
	本剧以校园教育为背景，采用黑色幽默、夸张表演与象征性影像，讽刺教育形式主义、机械服从及评价失真等现象。作品不追求写实再现，而是通过高度风格化的舞台语言，将日常校园生活中习以为常的荒诞加以放大，使观众在笑声与不安之间重新审视教育的真实处境。全剧以三段式结构展开，分别对应会议、兴趣班与视察三个典型场景，层层递进，最终指向对教育评价体系与独立人格的追问。
	
	\section{【创作主题】}
	
	作品通过学校会议、周六兴趣班和领导视察三个场景，展现管理者追求表面成果、教师难以表达异议、学生熟练背诵标准答案的荒诞现实，揭示教育形式与实际效果之间的矛盾。会议桌上，发言者照本宣科，数据被反复修饰，问题被巧妙回避；兴趣班里，孩子们以整齐划一的动作完成"个性发展"的表演，兴趣沦为任务；视察现场，一切井然有序，却在灯光与口号的缝隙中暴露出空洞。剧作试图追问：当服从被当作美德、当标准答案取代独立思考，教育究竟在培养什么样的人？形式主义的繁盛之下，真实的教育目标正在被悄然掏空。
	
	\section{【表现形式】}
	
	舞台采用统一的群体动作、重复音效、明暗灯光和影像投影，营造压抑而荒诞的氛围。演员以近乎机械的节奏完成动作，台词重复、错位，形成黑色幽默的张力；音效以钟声、掌声、翻页声等校园符号反复穿插，强化规训感。灯光在冷白与暖黄之间切换，暗示表演与真实的错位；影像投影呈现表格、标语与分数，构成对舞台事件的无声注解。开场与结尾的鲨鱼意象形成呼应，强化作品的象征意义——那潜伏于平静水面之下的巨大阴影，正是形式主义对教育本真的吞噬。
	
	\section{【预期效果】}
	
	通过讽刺与反转，引导观众反思教育评价体系、独立思考的价值，以及形式主义对真实教育目标的影响。作品希望观众在离场之后，仍能对身边习以为常的"标准流程"保持警觉，重新掂量服从与创造、表面与实质之间的分量，进而思考每一个教育参与者应承担的责任。
	
\end{document}